\section{Derivatives}

The derivative describes how a function changes locally. If
\[
y=f(x),
\]
then the derivative of \(f\) with respect to \(x\) is written as
\[
f'(x)
\qquad\text{or}\qquad
\frac{df}{dx}.
\]

\begin{definitionbox}
Let \(f\) be a real-valued function defined near \(a\). If the limit
\[
\lim_{h\to 0}\frac{f(a+h)-f(a)}{h}
\]
exists, then it is called the derivative of \(f\) at \(a\) and is denoted by \(f'(a)\).
\end{definitionbox}

For example, if
\[
f(x)=x^2,
\]
then
\[
f'(x)=2x.
\]
At \(x=1\), we get \(f'(1)=2\). At \(x=3\), we get \(f'(3)=6\). The function changes more rapidly near \(x=3\) than near \(x=1\).

If
\[
f(x)=3x+2,
\]
then
\[
f'(x)=3.
\]
The derivative is constant because the function is linear and has the same slope everywhere. If \(f(x)=C\), where \(C\) is constant, then \(f'(x)=0\).

The derivative also has a geometric interpretation. It is the slope of the tangent line to the graph of the function. If the derivative is positive, the function is increasing. If the derivative is negative, the function is decreasing. If the derivative is zero, the function is locally flat.

Some basic differentiation rules are
\[
\frac{d}{dx}x^n=nx^{n-1},
\qquad
\frac{d}{dx}(f(x)+g(x))=f'(x)+g'(x),
\qquad
\frac{d}{dx}(cf(x))=cf'(x),
\]
where \(c\) is a constant.

\begin{examplebox}
\[
\frac{d}{dx}(3x^2+5x-7)=6x+5.
\]
\end{examplebox}

In this chapter, the important point is purely mathematical: a derivative measures local change.
